\subsection{Assumptions}\label{sec:assumptions}

The analysis rests on three standing conditions.  They separate the regularity needed for uniform control, the structural source of policy-parameter dependence, and the smoothness needed for the nonlinear CPT expansion.  Conditions that enter only a specific comparison or corollary---such as reference-isolated coupling, stationary-set geometry, or the minimax oracle experiment---are introduced with the corresponding result rather than treated as global assumptions.

\begin{assumption}[Compact policy and trajectory domain]\label{ass:trajectory}
The parameter set $\Theta\subset\R^d$ is nonempty, compact, and convex, and is contained in an open set on which the policy is defined.  The horizon $h$ and asset count $N$ are finite.  There exist constants $Q_f,B_X,B_r,B_v<\infty$ such that, uniformly over every analyzed true model $P_k$ and frozen simulator $\widetilde P_k$,
\begin{equation}\label{eq:primitive-bounds}
 \sup_{i,u}\norm{q_{i,u}}_2\le Q_f,
 \qquad
 0<X_u\le B_X,
 \qquad
 0\le r_u\le B_r,
 \qquad
 0\le v_\pm(Y_k)\le B_v.
\end{equation}
\end{assumption}

Assumption~\ref{ass:trajectory} supplies a common analysis domain on which score moments, objective values, and smoothness constants can be controlled uniformly across episodes.  The quantities $\Theta,h,N$, and the feature norms are design-controlled, while the trajectory bounds are directly diagnosable from simulated or replayed paths.  If only a local or high-probability bounded region is available, the fixed-target arguments below still apply on that region; the online statements then inherit local or high-probability constants instead of the episode-uniform constants used here.

\begin{assumption}[Policy-mediated likelihood structure]\label{ass:likelihood}
Conditional on the observed history and executed action, the market transition kernel and the history-measurable feasibility map $\mathsf M_u(\cdot\mid s_u)$ have no explicit dependence on $\theta$.  The latent Dirichlet proposals $\widetilde\omega_u$ are retained in the augmented trajectory, and each frozen simulator is held fixed as $\theta$ varies.  Hence the policy density is the only explicit source of $\theta$-dependence in the conditional augmented-trajectory law.
\end{assumption}

This condition isolates the parameter dependence that the trajectory score in \eqref{eq:trajectory-score} differentiates.  It is auditable from the simulator interface: once the observed history and executed action are fixed, neither the market kernel nor $\mathsf M_u(\cdot\mid s_u)$ should take $\theta$ as an additional input.  If a model instead contains direct environment dependence on $\theta$, the corresponding likelihood or pathwise derivative must be added to the score identity; the estimator developed here then targets the policy-mediated component unless that additional term is included.

\begin{assumption}[Smooth regularized CPT weighting]\label{ass:cpt-regularity}
The preference parameters in \eqref{eq:value-functions}--\eqref{eq:normalized-weight} are fixed.  The normalized weighting functions $w_\pm^\eps$ are strictly increasing and belong to $C^3([0,1])$, with
\begin{equation}\label{eq:weight-derivative-bounds}
 C_j
 :=\max_{\sigma\in\{+,-\}}\sup_{p\in[0,1]}
 \abs{(w_\sigma^\eps)^{(j)}(p)}
 <\infty,
 \qquad j=1,2,3.
\end{equation}
\end{assumption}

Assumption~\ref{ass:cpt-regularity} is a property of the chosen regularized weighting functions and can therefore be checked analytically.  The first derivative enters the score-gradient identity, the second controls objective smoothness, and the third controls the second-order remainder used by the full/half debiasing construction.  Weaker integrability conditions can suffice for a gradient identity, but without uniform endpoint control the finite-sample expansion and finite-variance randomized correction require additional tail or margin conditions.

Under Assumptions \ref{ass:trajectory}--\ref{ass:cpt-regularity}, Appendix \ref{app:notation} establishes finite uniform score moments and constants $B_g,L<\infty$ such that, for both true and frozen-simulator objectives,
\begin{equation}\label{eq:gradient-smoothness-summary}
 \sup_{k,\theta\in\Theta}\norm{g_k(\theta)}_2\le B_g,
 \qquad
 \norm{g_k(\theta)-g_k(\theta')}_2
 \le L\norm{\theta-\theta'}_2.
\end{equation}
In particular, the simulator-gradient discrepancy $\beta_k$ and objective-variation budget $\cV_K$ in \eqref{eq:simulator-gradient-error}--\eqref{eq:objective-variation} are finite for every finite experiment.